\subsection{Mô hình Thread Pool}

Trong phương pháp lập trình đa luồng đơn sơ, chiến lược định tuyến dựa trên nguyên tắc khởi tạo một luồng thực thi (Thread) mới tương ứng với mỗi yêu cầu đầu vào (thread-per-request) bộc lộ giới hạn nghiêm trọng về thông lượng khi hệ thống tiếp nhận tải trọng lớn. Quá trình tạo lập và hủy phân bổ luồng liên tục yêu cầu chi phí tài nguyên đáng kể về chu kỳ CPU (CPU Cycles) và dung lượng vùng nhớ cục bộ (Thread Stack). Thêm vào đó, việc duy trì đồng thời hàng nghìn luồng độc lập dẫn đến suy giảm hiệu năng tổng thể của hệ điều hành, xuất phát từ hao phí chuyển đổi ngữ cảnh (Context Switching) và sự cạnh tranh bộ định thời (Scheduling Contention).

Mô hình Thread Pool là một mẫu thiết kế phần mềm (Design Pattern) được sử dụng nhằm hệ thống hóa vòng đời quản trị tài nguyên tính toán song song. Khác với cơ chế phân rã động, kiến trúc này khởi tạo sẵn một lượng luồng làm việc (Worker Threads) xác định và đặt chúng vào trạng thái chờ tĩnh (Idle). Các luồng thực thi này được tái sử dụng xoay vòng để giải quyết chuỗi tác vụ đẩy vào hệ thống thông qua một bộ đệm dữ liệu trung gian, phổ biến nhất là Hàng đợi tác vụ (Task Queue).

Việc ứng dụng cấu trúc Thread Pool mang lại các cơ chế tối ưu hóa hệ thống cốt lõi sau:
\begin{itemize}[label={-}]
    \item Tái sử dụng vùng nhớ cấp phát (Resource Reuse): Triệt tiêu chi phí tương tác Kernel (Kernel interaction) phát sinh từ tác vụ sinh và hủy luồng liên tục. Luồng làm việc tự động thu nhận tác vụ tiếp theo từ hàng đợi ngay sau khi hoàn tất phiên xử lý hiện tại, duy trì trạng thái hoạt động dài hạn.
    \item Giới hạn khả năng tương tranh (Concurrency Bounding): Thông qua việc cấu hình giới hạn tĩnh đối với số lượng luồng làm việc (thông thường được tham chiếu theo cấu trúc lõi vật lý của vi xử lý), kiến trúc thiết lập cơ chế giới hạn tài nguyên hệ thống (Resource Limiting). Yếu tố này củng cố khả năng phòng vệ của hệ thống trước tình trạng cạn kiệt bộ nhớ (OOM - Out of Memory) do xử lý quá tải.
    \item Cấu trúc điều phối đồng bộ chặt chẽ: Mô hình tích hợp các cơ chế đồng bộ hóa luồng tiêu chuẩn như Khóa loại trừ lẫn nhau (Mutex Lock) và Biến điều kiện (Condition Variable). Sự kết hợp này bảo vệ tính nguyên tử (Atomicity) khi các luồng cùng thực thi thao tác cấu trúc trên Hàng đợi tác vụ, đồng thời cho phép đánh thức luồng tĩnh một cách chính xác.
\end{itemize}

Trong sơ đồ vận hành của Secure Chat Server, Thread Pool được triển khai tuân theo mẫu thiết kế Producer-Consumer. Luồng thực thi mạng (Event Loop kết hợp API \texttt{epoll}) đảm nhiệm vai trò Producer, liên tục chuyển đổi các luồng byte thu nhận thành các khối tác vụ và nạp vào Hàng đợi. Tập hợp Worker Threads đảm nhiệm vai trò Consumer, tiến hành trích xuất tác vụ để thực hiện các tiến trình giải mã, tính toán xác thực, truy vấn dữ liệu và định tuyến lại mạng lưới. Sự phân tách trách nhiệm (Decoupling) này bảo đảm luồng I/O mạng luôn ở trạng thái đáp ứng (Responsive), hoàn toàn cách ly khỏi độ trễ do các thao tác sử dụng nhiều CPU (CPU-bound) gây ra.
